% Latex template: https://github.com/mqTeXUsers/Macquarie-University-Beamer-Theme

% Slide Masters:

% Title
% Text
% 2 column
% Full-image
% Bibliography
% Closing
\PassOptionsToPackage{table}{xcolor}    % https://tex.stackexchange.com/a/5365/5483
\documentclass[aspectratio=169, 12pt]{beamer} % Aspect ratio
% https://tex.stackexchange.com/a/14339/5483 
% Possible values: 1610, 169, 149, 54, 43 and 32.
% 169 = 16:9

\usetheme{macquarie}

\usepackage[english]{babel}       % Set language
\usepackage[utf8x]{inputenc}      % Set encoding
\usepackage{colortbl}
\mode<presentation>           % Set options
{
  \usetheme{default}          % Set theme
  \usecolortheme{default}         % Set colors
  \usefonttheme{default}          % Set font theme
  \setbeamertemplate{caption}[numbered] % Set caption to be numbered
}

% Uncomment this to have the outline at the beginning of each section highlighted.
%\AtBeginSection[]
%{
%  \begin{frame}{Outline}
%    \tableofcontents[currentsection]
%  \end{frame}
%}

\usepackage{graphicx}         % For including figures
\usepackage{booktabs}         % For table rules
\usepackage{hyperref}         % For cross-referencing


\setbeamertemplate{footline}[frame number] 

\title{Title for a minimal beamer presentation} % Presentation title
\author{Author One}               % Presentation author
\institute{Name of department/affiliation}         % Author affiliation
\date{\today}                 % Today's date  
\begin{document}

% Title page
% This page includes the informations defined earlier including title, author/s, affiliation/s and the date
% \begin{frame}[noframenumbering]

\maketitle

  
% \end{frame}

% Outline
% This page includes the outline (Table of content) of the presentation. All sections and subsections will appear in the outline by default.
\begin{frame}{Outline}
  \tableofcontents
\end{frame}

% The following is the most frequently used slide types in beamer
% The slide structure is as follows:
%
%\begin{frame}{<slide-title>}
% <content>
%\end{frame}

\section{Section One}

\begin{frame}{Slide with bullet points}
  This is a bullet list of two points:
    \begin{itemize}
    \item Point one
        \item Point two
  \end{itemize}
\end{frame}

\begin{frame}{Slide with two columns}
  \begin{columns}
    \column{.5\textwidth}
        Text goes in first column.
        
        \column{.5\textwidth}
        Text goes in second column
  \end{columns}
\end{frame}

\section{Section Two}

\begin{frame}{Slide with table}
  % Please add the following required packages to your document preamble:
% \usepackage{booktabs}

% \definecolor{mqRed}{HTML}{A6192E}
% \definecolor{mqDeepRed}{HTML}{76232F}
% \definecolor{mqBrightRed}{HTML}{D6001C}
% \definecolor{mqMagenta}{HTML}{C6007E}
% \definecolor{mqPurple}{HTML}{80225F}
% \definecolor{mqCharcoal}{HTML}{373A36}
% \definecolor{mqSand}{HTML}{D6D2C4}
% \definecolor{mqBlack}{HTML}{000000}
% \definecolor{mqWhite}{HTML}{FFFFFF}
\begin{table}[H]
\centering

%\rowcolors{1}{mqSand}{mqSand!50}
\newcolumntype{a}{>{\columncolor{mqPurple}}l}
\newcolumntype{b}{>{\columncolor{mqSand}}c}

\begin{tabular}{abb}

\rowcolor{mqCharcoal}
\multicolumn{3}{l}{\textcolor{white}{\textbf{Caption for table one}}}\\

\rowcolor{mqBrightRed} 

\textcolor{white}{\textbf{Heading 1}} & \textcolor{white}{\textbf{Heading 2}} & \textcolor{white}{\textbf{Heading 3}} \\ 
\textcolor{white}{Row 1} & 0.1 & .01 \\
\textcolor{white}{Row 2} & 0.2 & .02 \\
\textcolor{white}{Row 3} & 0.3 & 0.03 \\
\textcolor{white}{Row 4} & 0.4 & 0.04 \\ 


% \toprule
% \textbf{Heading 1} & \textbf{Heading 2} & \textbf{Heading 3} \\ \midrule
% Row1 & 0.1 & .01 \\
% Row2 & 0.2 & .02 \\
% Row3 & 0.3 & 0.03 \\
% Row4 & 0.4 & 0.04 \\ \bottomrule
\end{tabular}
\end{table}
\end{frame}

\begin{frame}{Slide with different table}
  % Please add the following required packages to your document preamble:
% \usepackage{booktabs}

% \definecolor{mqRed}{HTML}{A6192E}
% \definecolor{mqDeepRed}{HTML}{76232F}
% \definecolor{mqBrightRed}{HTML}{D6001C}
% \definecolor{mqMagenta}{HTML}{C6007E}
% \definecolor{mqPurple}{HTML}{80225F}
% \definecolor{mqCharcoal}{HTML}{373A36}
% \definecolor{mqSand}{HTML}{D6D2C4}
% \definecolor{mqBlack}{HTML}{000000}
% \definecolor{mqWhite}{HTML}{FFFFFF}
\begin{table}[H]
\centering
\caption{Caption for table one}
\label{tab:table1}
\rowcolors{1}{mqSand}{mqSand!50}
\begin{tabular}{@{}lcc@{}}
\rowcolor{mqBrightRed}

\textcolor{white}{\textbf{Heading 1}} & \textcolor{white}{\textbf{Heading 2}} & \textcolor{white}{\textbf{Heading 3}} \\ 
Row 1 & 0.1 & .01 \\
Row 2 & 0.2 & .02 \\
Row 3 & 0.3 & 0.03 \\
Row 4 & 0.4 & 0.04 \\ 
\rowcolor{mqDeepRed}
asdf & qwer & asdf

% \toprule
% \textbf{Heading 1} & \textbf{Heading 2} & \textbf{Heading 3} \\ \midrule
% Row1 & 0.1 & .01 \\
% Row2 & 0.2 & .02 \\
% Row3 & 0.3 & 0.03 \\
% Row4 & 0.4 & 0.04 \\ \bottomrule
\end{tabular}
\end{table}
\end{frame}

\begin{frame}{Slide with figure}
  \begin{figure}[H]
    \centering
        \includegraphics[height=.7\textheight]{branding/MQ_INT_VER_RGB_POS.pdf}
        \caption{Caption for figure one. MQ logo vertical.}
        \label{fig:figure1}
  \end{figure}
\end{frame}

\begin{frame}{Slide with references}
  This is to reference a figure (Figure \ref{fig:figure1})\\
    This it to reference a table (Table \ref{tab:table1})\\
    This is to cite an article \cite{Sobotkova2015-lq, Nosek2018-yv}\\
    This is to add an article to the references without mentioning in the text \nocite{Borgman2015-rh, Center_for_Open_Science2018-pn}\\
\end{frame}
\section{References}

% Adding the option 'allowframebreaks' allows the contents of the slide to be expanded in more than one slide.
% The "1" comes from the outer theme"
\begin{frame}[allowframebreaks]{References}
  \tiny
  \bibliography{references}
  \bibliographystyle{apalike}
\end{frame}


\begin{frame}[standout]
Thank you!
\end{frame}



\end{document}
